% Generated draft: requires mathematical and source review.
\documentclass[11pt]{article}
\usepackage[margin=25mm]{geometry}
\usepackage{fontspec}
\setmainfont{Latin Modern Roman}
\usepackage{amsmath,amssymb,mathtools,unicode-math}
\setmathfont{Latin Modern Math}
\usepackage{enumitem,needspace,xurl,hyperref}
\setlength{\parindent}{0pt}
\setlength{\parskip}{5pt}
\setlength{\emergencystretch}{3em}
\newcommand{\N}{\mathbb{N}}
\newcommand{\Z}{\mathbb{Z}}
\newcommand{\Q}{\mathbb{Q}}
\newcommand{\R}{\mathbb{R}}
\newcommand{\C}{\mathbb{C}}
\newcommand{\F}{\mathbb{F}}
\newcommand{\A}{\mathbb{A}}
\newcommand{\PP}{\mathbb{P}}
\newcommand{\E}{\mathbb{E}}
\newcommand{\eps}{\varepsilon}
\newcommand{\dd}{\,\mathrm{d}}
\begin{document}
\section*{UnsolvedMath ID: 10400077}
\textbf{Source problem code:} catalogue entry
\subsection{Definitions and mathematical statement}
Let $\Sigma$ be the surface and $G$ the group in the intended twisted Wess--Zumino--Witten model. A flat principal $G$-connection is specified up to gauge by its holonomy data. After choosing the level, insertion labels, marked points, and twisting data, conformal blocks form the source's vector-space/bundle family.

Compute the holonomy of this family. In particular, for each flat $G$-connection on $\Sigma$, determine the induced surface-braid action on the corresponding conformal blocks. For $n$ unordered distinct points, the usual surface braid group is
\[
B_n(\Sigma)=\pi_1(\operatorname{Conf}_n(\Sigma)/S_n),
\]
while ordered points give the pure braid group. If the relevant connection on the conformal-block bundle is only projectively flat, its holonomy is a representation into $\operatorname{PGL}(V)$; a linear lift requires additional normalization/central-extension data. The source does not specify this distinction or the labels needed for full versus pure/colored braids.

The extract also names $\mathcal A_n(\Sigma)$, the algebra of chord diagrams on $n$ strings with horizontal chords on $\Sigma$. It supplies neither the surface relations nor a precise monodromy/completion comparison; this is retained as context, not an added basis-computation task.

The surface category, group, WZW level and insertion representations, twist, parameter space, projective connection, and behavior of flat-connection data under braid transport are all missing. These must be fixed to identify the holonomy representation being requested; an arbitrary untwisted braid representation is not a substitute for the twisted, each-flat-connection question.
\subsection{Short English statement}
Compute twisted-WZW conformal-block holonomy and the surface braid action for each flat $G$-connection. Fix the model and labeling data and distinguish projective from linear and full from pure/colored braid representations.
\subsection{Sources}
\begin{itemize}
\item Ohtsuki, \emph{Problems on invariants of knots and 3-manifolds} (2002), Problem 3.27, PDF page 72, attributed to T. Kohno, \url{https://msp.org/gtm/2002/04/gtm-2002-04-024s.pdf}; accessed 2026-07-29, PDF-font-structure extraction.
\item Local registry \texttt{data/unsolvedmath/problems.json}, permanent UnsolvedMath ID $10400077$, source code catalogue entry, \emph{Problem 3.27}. Dataset attribution: Ulam AI and the UnsolvedMath Contributors.
\end{itemize}
\end{document}
